(1) 三维自由粒子的单粒子态密度为 \(D(\varepsilon)=V(2m)^{3/2}\varepsilon^{1/2}/(4\pi^2\hbar^3)\)。把 \((z^{-1}e^{\beta\varepsilon}-1)^{-1}\) 展为 \(\sum_{\ell\ge1}z^\ell e^{-\ell\beta\varepsilon}\)，逐项用 \(\Gamma(3/2)\) 与 \(\Gamma(5/2)\) 积分，得到 \(N_{\rm ex}=V\lambda_T^{-3}g_{3/2}(z)\)、\(U=(3/2)Vk_BT\lambda_T^{-3}g_{5/2}(z)\)。零能基态单独计数 \(N_0=N-N_{\rm ex}\)，不能用连续态密度替代其宏观占据。

(2) 对自由玻色气体 \(\mu\le0\)，故 \(z=e^{\beta\mu}\le1\)，激发态容纳量上界为 \(V\lambda_T^{-3}\zeta(3/2)\)。当此上界等于总粒子数 \(nV\) 时 \(T_c=\frac{2\pi\hbar^2}{mk_B}[n/\zeta(3/2)]^{2/3}\)。低于此温度 \(z=1,\mu=0\)，于是 \(N_0/N=1-(T/T_c)^{3/2}\)。二维自由态密度为常数，\(g_1(z)\) 在 \(z\uparrow1\) 发散，因而均匀二维理想气体在有限温度下无同类热力学凝聚。

(3) 低温 \(U=(3/2)Nk_BT\,\zeta(5/2)\zeta(3/2)^{-1}(T/T_c)^{3/2}\)，固定 \(N,V\) 求导得 \(C_V/Nk_B=(15/4)[\zeta(5/2)/\zeta(3/2)](T/T_c)^{3/2}\)。在正常相，\(N=V\lambda_T^{-3}g_{3/2}(z)\)，固定粒子数微分并用 \(z\,g_\nu'(z)=g_{\nu-1}(z)\) 得 \(d\ln z/d\ln T=-(3/2)g_{3/2}(z)/g_{1/2}(z)\)。于是 \(C_V/Nk_B=(15/4)g_{5/2}/g_{3/2}-(9/4)g_{3/2}/g_{1/2}\)。当 \(z\uparrow1\)，\(g_{1/2}\to\infty\)，两侧同趋 \((15/4)\zeta(5/2)/\zeta(3/2)\)，热容连续但非解析。

(4) 巨配分函数的体积导数给 \(P=k_BT\lambda_T^{-3}g_{5/2}(z)\)。与第一问能量公式比较，精确得到 \(PV=2U/3\)，亦可从动量通量或非相对论尺度齐次性独立验证。凝聚态基态动能为零，在热力学极限的压力贡献 \(O(1/V)\) 消失。\(T\to0\) 时 \(U\propto T^{5/2}\)、\(P\propto T^{5/2}\)、\(C_V\propto T^{3/2}\) 均趋零；相互作用存在时则需引入不同的低能声子物理，不能外推理想气体结果。
